\begin{frontmatter}

\title{Frequent Losers as Cover Bidders in Public Procurement\texorpdfstring{\footnote{This version: March 2026.}}{}}

\author[insper]{Darcio Genicolo-Martins\corref{cor1}}
\ead{darciogm1@insper.edu.br}
\author[insper]{Paulo Furquim de Azevedo}
\affiliation[insper]{organization={INSPER},
  city={S\~ao Paulo}, country={Brazil}}
\cortext[cor1]{Corresponding author.}

\begin{abstract}
Bid-rigging cartels must simulate competition to avoid detection,
often by deploying firms that submit deliberately losing bids. We
propose a firm-level screen based on this footprint: \emph{frequent
losers}---firms that never win yet participate in abnormally many
tenders. Using 4.5 million tender-items from Brazilian public
procurement (2009--2019), we show that tenders with frequent-loser
participation exhibit 3.6--8.0\% higher prices across five
independent estimators, and that 7.1\% of flagged firms
co-participate with competition-authority-convicted cartelists---3.5
times the expected rate. The screen achieves AUC~$= 0.94$ against
known cartels using only participation records, outperforming
bid-level screens (AUC~$= 0.79$). Convergent evidence points to a
cover-bidding mechanism: the price effect concentrates in competitive
markets, bid-coordination tests reject exchangeability and
conditional independence of frequent-loser residuals, and structural
estimation identifies coordinated cover-bid placement. The screen
requires no bid-level data, making it operable where existing
forensic tools cannot reach.
\end{abstract}

\end{frontmatter}

\noindent\textbf{Keywords:} bid rigging, cover bidding, public procurement, cartel screening, frequent losers,\\ instrumental variables

\noindent\textbf{JEL No.:} D44, D73, H57, K21, L41

\bigskip
